\section{Interfaces, Data, and Ecosystem}
\label{sec:interfaces}

\subsection{The Facade: A Single Task-Level Entry}

All capabilities converge on a single \texttt{Facade} class (ADR 0014), from
which the CLI and the MCP server are derived from the same metadata source:
beyond input validation, the task-level API model also exposes parameter
metadata for the current task context (conditional value domains and
validators share one rule definition), so that GUIs, the CLI, and MCP can
generate widgets and range hints---external consumers keep no local copies of
parameter rules. Task-level methods come in three tiers:

\begin{itemize}
  \item \textbf{Primary tasks}: \texttt{design\_orbit} (mission orbit design:
        CR3BP initial guess $\to$ ephemeris correction $\to$ high-fidelity
        propagation), \texttt{control\_orbit} (station-keeping Monte Carlo
        simulation), \texttt{transfer\_design} (HMN/LGA/WSB/low-thrust),
        \texttt{orbit\_propagation} (orbit propagation), and
        \texttt{spacetime\_transform} (synodic$\leftrightarrow$J2000 and
        GCRS$\leftrightarrow$EBCRS conversions);
  \item \textbf{Secondary subtasks}: \texttt{orbit\_family\_generation}
        (continuation generation of eight families, with family records
        automatically archived), among others;
  \item \textbf{Orbit catalog}: 7 catalog methods (below).
\end{itemize}

\begin{lstlisting}[caption={Designing an Earth--Moon L2 Halo orbit (examples/main\_design.py)},label=lst:quickstart]
from e2m2e.api import Facade

facade = Facade()

result = facade.design_orbit(
    orbit_type="Halo",
    collinear_point=2,
    amplitude=30000.0,          # km
    epoch=[2024, 1, 1, 0, 0, 0.0],
    duration=365.25 * 86400.0,  # s
)

print(result.orbit_type)
print(result.initial_state)
\end{lstlisting}

\subsection{MCP: An Agent-Native Interface}

e2m2e can act as an MCP server (\texttt{e2m2e mcp-serve}, stdio transport,
listening on no port), exposing 13 task-level tools to LLM agents: 6 task
methods (the 5 primary ones plus family generation) and 7 orbit-catalog
tools. The tool list is purely derived from Facade method metadata---nothing
else to maintain. All responses use a unified envelope
\texttt{\{status, data, error, meta\}}: errors are expressed as error codes
(\texttt{INVALID\_PARAMS}, \texttt{RECORD\_NOT\_FOUND},
\texttt{PROPAGATION\_FAILED}, etc.) without leaking tracebacks, and numerical
tools report soft-failure semantics via \texttt{data.status} (converged /
diverged / stagnated / max\_iterations, ADR 0020)---an agent can treat
``not converged'' as normal information rather than an exception.

Setup is a single command registered in the MCP client; it is recommended to
pin \texttt{cwd} to the directory containing \texttt{kernels/} (SPICE
kernels) and \texttt{catalog/} (the orbit catalog), or to give absolute paths
via \texttt{SPICE\_KERNEL\_DIR} / \texttt{E2M2E\_CATALOG\_DIR}. Once
configured, the server can be driven in natural language, e.g.:

\begin{quote}\itshape
Design an L2 southern NRHO with 3000 km perilune altitude, then run 100
Monte Carlo station-keeping simulations on it and tag the results
``candidate''.
\end{quote}

\subsection{The Orbit Catalog}

Computation artifacts are automatically archived (ADR 0031): a record is JSON
metadata + NPZ array segments (with a schema version), with SQLite only as a
derived index. Productive tools return a \texttt{record\_id} on success,
which chains across tools (e.g. \texttt{control\_orbit} can reference the
output of \texttt{design\_orbit} via \texttt{input\_record\_id}). Seven
catalog tools cover the full lifecycle: \texttt{catalog\_query}
(multi-dimensional filtering by family / libration point / Jacobi range /
amplitude / tag / convergence status), \texttt{catalog\_get},
\texttt{catalog\_delete}, \texttt{catalog\_tag} (teaching annotations),
\texttt{catalog\_promote} (promoting a family member to a standalone record),
\texttt{catalog\_export} (packaged export), and \texttt{catalog\_sweep}
(batch generation over parameter grids).

\subsection{Installation and Deployment}

\begin{lstlisting}[style=e2bash,caption={Installation and SPICE kernel setup},label=lst:install]
uv pip install e2m2e            # release wheels embed CSPICE, ready out of the box
uv pip install "e2m2e[mcp]"     # MCP server

# development from source (requires the Rust toolchain):
git clone https://github.com/cislunarspace/e2m2e.git
cd e2m2e
make dev      # sync deps + fetch CSPICE package & SPICE kernels + build the Rust extension
\end{lstlisting}

Release wheels cover Windows x86\_64, Linux x86\_64, and Linux aarch64
(Kunpeng / Phytium / Raspberry Pi), with statically linked CSPICE embedded,
so the Rust fast paths (STM propagation, shooting, third-body gravity, etc.)
work out of the box; when the extension is unavailable, an explicit error is
raised rather than silently degrading. The SPICE kernels are packaged as the
GitHub Release asset \texttt{kernels-v1} (downloaded automatically by
\texttt{make setup}), including DE planetary ephemerides, Earth rotation,
Moon attitude, leap seconds, and planetary constants kernels; the r2s2
spacetime transformation additionally requires an ephemeris containing the
TT$-$TDB time ephemeris. The official source is NASA NAIF.

The \texttt{examples/} directory provides six runnable examples: end-to-end
orbit design, station-keeping Monte Carlo, high-fidelity propagation, a
Lambert transfer, a four-layer comparison of lunar orbits (LLO/ELFO/DRO/Halo,
under one minute in total), and orbit-family generation.
